\section{Messung der Streuspektren}
\subsection{Theoretische Grundlagen}
Aus der relativistischen Compton-Formel !!!referenz compton formel in abschnitt 1!!! folgt, dass der Vollenergiepeak des Primärstrahls bei zunehmendem Streuwinkel sich zu niedrigeren Energien verschiebt. Gleichzeitig ensteht ein kontinuierliches Spektrum von Compton-Photonen, die durch die Streuung an Elektronen entstehen: das Compton-Kontinum. Das Spektrum enthält somit sowohl den Vollenergiepeak als auch das Compton-Kontinuum, das von der Energie des Primärstrahls bis zum Compton-Edge geht.\\
Durch die Messung der Spektrum für unterschiedliche Streuwinkel $\theta$ kann die Winkelabhängigkeit der gestreuuten Photonenenergie experimentel überprüft werden.\\

( Es ist zu erwarten, dass mit zunehmendem Streuwinkel die Intensität des Vollenergiepeaks abnimmt, während das Compton-Kontinuum an Bedeutung gewinnt. Insbesondere bei größeren Streuwinkeln wird der Vollenergiepeak deutlich abgeschwächt, da mehr Photonen Energie an die Elektronen abgeben und somit im Compton-Kontinuum landen.)
\subsection{Durchführung}
Zur Erzeugung eines definierten Einzelstreuprozesses, wurde eine 1 mm dicke Aluminiumfolie auf den kollimierten Gammastrahl gelegt. Die gestreuten Photonen wurden in 5°-Schritten unter einem Winkel von 35° bis 120° detektiert. Für diese Messung wurde ein 3 cm großer Kollimator verwendet, um eine hohe Messgeschwindigkeit zu gewährleisten. Da die Streueffizienz einzelner Teilchen im Target geringer ist, erhöht ein größerer Kollimator die Empfindlichkeit des Detektors, ohne die Winkelauflösung wesentlich zu beeinträchtigen.
\paragraph{Untergrundmessung}
Ohne Target gelangen auch gestreute Photonen im Detektor. Dies kann darauf zurückgeführt werden, dass durch die Geometrie des Detektors, ein Teil der Primärstrahlung an den Bleiwänden gestreut wird und in fast allen Winkeln ausreteten kann. Die Photonen können auch an Detektormaterial selbst gestreut werden, sowie an Luftmolekülen und Aufbauten. \\
Diese Beiträge sind winkelabhängig und können nicht vernachlässigt werden. Dies führt dazu, dass für jede Detektorposition eine neue, separate Untergrundmessung durchgeführt werden muss, damit die Spektren korrekt interpretiert werden können. \\
Daher wird in der Praxis für jede Position des Detektors eine Untergrundmessung ohne Target durchgeführt.\\
Die Differenz $N_{new}(\theta)=N_{mit Target}-N_{ohne Target}$ gibt das reine Streuspektrum des Aluminumtargets an. 